% !TeX program = pdflatex
% Lectura opcional: no forma parte del resumen de contenidos ya vistos.
\documentclass[10pt,a4paper]{article}
\makeatletter
\def\input@path{{../semana_03/}}
\makeatother
\usepackage{estilo_apuntes}
\geometry{a4paper,left=14mm,right=14mm,top=22mm,bottom=16mm,headheight=32pt}
\usepackage{multicol}
\setlength{\columnsep}{7mm}
\newcommand{\R}{\mathbb{R}}
\renewenvironment{definicion}[1][]{\par\medskip\noindent{\color{Azul}\bfseries Definición: #1}\par\nobreak\smallskip}{\par}
\encabezado{Lectura opcional}{Formalización de límites}
\pagestyle{fancy}
\begin{document}
\begin{center}
  {\Large\bfseries\color{Azul} Formalización de límites y continuidad}
\end{center}
\textbf{Para un estudio posterior.} Estas definiciones se reservan para cuando se introduzcan los cuantificadores y la formulación precisa de proximidad. No forman parte de los contenidos ya desarrollados en clase.

\textbf{Lectura de los símbolos.} La condición $|x-a|<\delta$ controla la distancia de la entrada a $a$; $|f(x)-L|<\varepsilon$ controla la distancia de la imagen a $L$. Primero se fija una tolerancia arbitraria $\varepsilon>0$ y luego se busca un $\delta>0$ que garantice la condición para todas las entradas indicadas. El valor de $\delta$ puede depender de $\varepsilon$, pero no de cada entrada $x$.
\begin{multicols}{2}
\begin{definicion}[límite y límites laterales]
Sea $f:D\to\R$ y sea $a$ un punto de acumulación de $D$: para todo $\delta>0$ existe $x\in D$ con $0<|x-a|<\delta$. Para $L\in\R$, se define
\[
 \lim_{x\to a}f(x)=L
\]
si para todo $\varepsilon>0$ existe $\delta>0$ tal que, para todo $x\in D$,
\[
 0<|x-a|<\delta\ \Longrightarrow\ |f(x)-L|<\varepsilon.
\]
Los límites laterales se definen con los mismos cuantificadores:
\begin{itemize}
  \item Por la izquierda se escribe
  \[
    \lim_{x\to a^-}f(x)=L,
  \]
  y se exige
  \[
    a-\delta<x<a\ \Longrightarrow\ |f(x)-L|<\varepsilon.
  \]
  \item Por la derecha se escribe
  \[
    \lim_{x\to a^+}f(x)=L,
  \]
  y se exige
  \[
    a<x<a+\delta\ \Longrightarrow\ |f(x)-L|<\varepsilon.
  \]
\end{itemize}
En cada caso, $a$ debe ser punto de acumulación del dominio por el lado correspondiente.
\end{definicion}

\begin{definicion}[límites infinitos y al infinito]
Sea $f:D\to\R$ y sea $a$ punto de acumulación de $D$. Se escribe
\[
  \lim_{x\to a}f(x)=+\infty
\]
si para todo $M>0$ existe $\delta>0$ tal que, para todo $x\in D$,
\[
 0<|x-a|<\delta\ \Longrightarrow\ f(x)>M.
\]
Para $\lim_{x\to a}f(x)=-\infty$, se sustituye la conclusión por $f(x)<-M$. Los límites laterales infinitos se definen restringiendo $x$ al lado correspondiente.

Si $D$ no está acotado superiormente y $L\in\R$, se escribe
\[
  \lim_{x\to+\infty}f(x)=L
\]
si para todo $\varepsilon>0$ existe $A>0$ tal que, para todo $x\in D$,
\[
 x>A\ \Longrightarrow\ |f(x)-L|<\varepsilon.
\]
En $\lim_{x\to-\infty}f(x)=L$ se exige que $D$ no esté acotado inferiormente y se sustituye $x>A$ por $x<-A$.

Se define
\[
  \lim_{x\to+\infty}f(x)=+\infty
\]
si para todo $M>0$ existe $A>0$ tal que, para todo $x\in D$,
\[
  x>A\ \Longrightarrow\ f(x)>M.
\]
Para límite $-\infty$ se exige $f(x)<-M$; en $x\to-\infty$ se usa $x<-A$, con el dominio no acotado inferiormente.
\end{definicion}

\begin{definicion}[continuidad]
Sea $f:D\to\R$ y $a\in D$. Se dice que $f$ es continua en $a$ si para todo $\varepsilon>0$ existe $\delta>0$ tal que, para todo $x\in D$,
\[
 |x-a|<\delta\ \Longrightarrow\ |f(x)-f(a)|<\varepsilon.
\]
Si $a$ es punto de acumulación de $D$, esta condición equivale a
\[
 \lim_{x\to a}f(x)=f(a).
\]
Una función es continua en un conjunto si es continua en cada uno de sus puntos. Para $f:[a,b]\to\R$, esto equivale a continuidad en $(a,b)$ y a
\[
 \lim_{x\to a^+}f(x)=f(a),\qquad
 \lim_{x\to b^-}f(x)=f(b).
\]
\end{definicion}
\end{multicols}
\end{document}
